\section{Preliminaries}\label{sec:preliminaries}

A sequence over a set \(X\) is an element of the free abelian monoid
\(\F(X)\). We write it multiplicatively, even when \(X\) is an
additively written group. Thus
\[
  S=x_1\cdot\ldots\cdot x_\ell,
  \qquad |S|=\ell.
\]
Subsequences are selected by positions, so equal values occurring in
different positions remain available separately. We write \(T\mid S\)
when \(T\) is a subsequence and \(ST^{-1}\) for the complementary
sequence. A sequence over a group is \emph{product one} if its terms can
be ordered so that their product is the identity.

For an abelian group \(B\), the Davenport constant \(\D(B)\) is the
least integer \(L\) such that every sequence of length \(L\) over
\(B\) has a nonempty zero-sum subsequence. Thus
\(\D(B)=\dsmall(B)+1\). If \(W\subseteq\mathbb Z\) is nonempty, the
weighted Davenport constant \(\D_W(B)\) is defined analogously, with a
coefficient from \(W\) assigned to each selected term. We write
\[
  \Dpm(B)=\D_{\{-1,1\}}(B).
\]

Throughout the paper, set
\[
  q_i=p^{\lambda_i},\qquad
  Q=|A|=\prod_{i=1}^r q_i,
  \qquad
  D_A=\D(A)=1+\sum_{i=1}^r(q_i-1).
\]
The formula for \(D_A\) is Olson's theorem
\cite{Olson1969}. In particular,
\begin{equation}\label{eq:basic-A-facts}
  D_A\le Q,
  \qquad \exp(A)\mid Q,
  \qquad D_A\ \text{and}\ Q\ \text{are odd}.
\end{equation}
The first inequality follows by expanding
\(Q=\prod_i(1+(q_i-1))\).

We represent elements of \(G=\Dih(A)\) as pairs
\((x,\varepsilon)\in A\times C_2\), with multiplication
\begin{equation}\label{eq:group-law}
  (x,\varepsilon)(y,\delta)
  =\bigl(x+(-1)^\varepsilon y,\,\varepsilon+\delta\bigr).
\end{equation}
Elements \((x,0)\) are called \emph{rotations}, and elements
\((x,1)\) are called \emph{reflections}.

\subsection{Realizing prescribed signs}

The next elementary lemma is the bridge between weighted sums in \(A\)
and product-one subsequences in \(G\).

\begin{lemma}[Realization of prescribed signs]\label{lem:sign-realization}
Suppose a selected sequence contains \(2k>0\) reflections and any
number of rotations. Assign signs from \(\{+1,-1\}\) to the terms so
that exactly \(k\) reflections receive each sign; the signs on the
rotations are arbitrary. Then the terms can be ordered so that their
contributions to the \(A\)-coordinate in \eqref{eq:group-law} have the
assigned signs. Consequently, if the assigned signed vector sum is
zero, the selected sequence is product one.

If a selected sequence contains no reflections, it is product one if
and only if the sum of its rotation vectors is zero.
\end{lemma}

\begin{proof}
Place the positively and negatively signed reflections alternately,
starting with a positive one. In any ordered product, a term contributes
its vector with sign \((-1)^t\), where \(t\) is the number of preceding
reflections. Hence the chosen reflection ordering realizes the prescribed
alternating signs.

The \(2k\) reflections determine \(2k+1\) gaps: one before the first
reflection, one after the last, and one between each consecutive pair.
A rotation inserted in a gap preceded by an even number of reflections
has sign \(+1\); in a gap preceded by an odd number it has sign \(-1\).
Since \(k\ge1\), gaps of both parities occur, and any number of rotations
can be inserted in either type of gap. This realizes all prescribed
rotation signs. The total number of reflections is even, so the
\(C_2\)-coordinate is zero; the \(A\)-coordinate is the prescribed
signed sum.
\end{proof}

\subsection{Prescribed-length weighted sums}

For a sequence \(T\) over a finite abelian group \(B\), a nonempty
weight set \(W\subseteq\mathbb Z\), and an integer \(n\le |T|\), let
\[
  \Sigma_n^W(T)
  =\left\{
      \sum_{j=1}^n w_jt_j:
      t_1,\ldots,t_n\text{ occur in distinct positions of }T,
      \ w_j\in W
    \right\}.
\]
For an integer \(n\), write \(nB=\{nx:x\in B\}\).
We use the following two corollaries of Grynkiewicz, Marchan, and Ordaz
\cite[Corollaries 1.2 and 1.3]{GrynkiewiczMarchanOrdaz2012}.

\begin{theorem}[Grynkiewicz--Marchan--Ordaz]\label{thm:GMO}
Let \(B\) be a finite abelian group, let \(W\subseteq\mathbb Z\) be
nonempty, let \(T\) be a sequence over \(B\), and let \(n\ge |B|\).
If
\[
  |T|\ge n+\D_W(B)-1,
\]
then
\begin{equation}\label{eq:GMO-existence}
  \Sigma_n^W(T)\cap nB\ne\varnothing.
\end{equation}
If, in addition, \(\gcd(W)=1\), then either
\begin{equation}\label{eq:GMO-full}
  \Sigma_n^W(T)=B,
\end{equation}
or there exist a proper subgroup \(H<B\), elements
\(\alpha,\beta\in B\), and a subsequence \(T'\mid T\) such that
\begin{equation}\label{eq:GMO-concentration}
  |T'|\ge |T|-|B/H|+2,
  \qquad
  \supp(T')\subseteq \alpha+H,
\end{equation}
and
\begin{equation}\label{eq:GMO-weight-coset}
  W x\subseteq\beta+H
  \quad\text{for every }x\in\supp(T').
\end{equation}
\end{theorem}

Two specializations will be used repeatedly. For \(W=\{1\}\), one has
\(\D_W(B)=\D(B)\). For \(W=\{-1,1\}\), condition
\eqref{eq:GMO-weight-coset} says that \(x\) and \(-x\) lie in the same
\(H\)-coset for every term of \(T'\). Thus \(2x\in H\); if \(B/H\)
has odd order, this implies \(x\in H\).

The other external input is the following consequence of
\cite[Theorem 1.1]{GodaraJoshiMazumdar2026}:
\begin{equation}\label{eq:small-Davenport-dihedral}
  \dsmall(\Dih(B))=\D(B)
  \qquad\text{for every finite abelian group }B.
\end{equation}

\subsection{The lower bound and identity padding}

\begin{proposition}\label{prop:lower-bound}
For \(G=\Dih(A)\),
\[
  \E(G)\ge 2Q+D_A.
\]
\end{proposition}

\begin{proof}
By \eqref{eq:small-Davenport-dihedral}, there is a product-one-free
sequence \(U\) over \(G\) of length \(D_A\). Append \(2Q-1\) copies
of the identity. Every subsequence of length \(2Q\) uses a nonempty
part of \(U\). If such a subsequence were product one, deleting the
identity terms from a product-one ordering would give a nonempty
product-one subsequence of \(U\), a contradiction.
\end{proof}

\begin{lemma}[Identity padding]\label{lem:identity-padding}
Let \(Y\) be a sequence over \(G\) of length \(2Q+D_A\). If at least
\(D_A\) terms of \(Y\) are the identity rotation, then \(Y\) contains
a product-one subsequence of length \(2Q\).
\end{lemma}

\begin{proof}
Let \(m_0\) be the number of identity rotations. From the sequence of
all nonidentity terms, successively remove nonempty product-one
subsequences until the remainder \(R\) is product-one-free. Let \(P\)
be the union of the removed subsequences. Concatenating one product-one
ordering for each removed block makes \(P\) product one. By
\eqref{eq:small-Davenport-dihedral}, \(|R|\le D_A\), and hence
\[
  |P|=(2Q+D_A-m_0)-|R|\ge 2Q-m_0.
\]
Since \(m_0\ge D_A\), also \(|P|\le 2Q\). Therefore at most \(m_0\)
unused identity terms are needed to extend \(P\) to length \(2Q\).
\end{proof}
